% -------------------------------------------------------------------------
% File: src/devops/github-actions/workflows-in-github-actions.tex
% Style: style v1.0 (see tooling/latex/style.sty)
% Build: pdflatex via latexmk; minted-aware with listings fallback.
% -------------------------------------------------------------------------
\documentclass[11pt]{article}
\usepackage{technical-installation}

% ---------- Encoding & layout ----------
\geometry{margin=1in}
\setstretch{1.1}

% ---------- Colors, links, boxes ----------
\definecolor{ink}{HTML}{111827}      % gray-900
\definecolor{soft}{HTML}{F9FAFB}     % gray-50
\definecolor{accent}{HTML}{2563EB}   % blue-600
\definecolor{ok}{HTML}{059669}       % emerald-600
\definecolor{warn}{HTML}{D97706}     % amber-600
\definecolor{bad}{HTML}{DC2626}      % red-600
% ---------- Silence first-run .toc warnings (latexmk treats them as missing input) ----------
\makeatletter
\patchcmd{\@starttoc}
  {\InputIfFileExists{\jobname.#1}{}{\typeout{No file \jobname.#1.}}}
  {\InputIfFileExists{\jobname.#1}{}{}}
  {}{}
\makeatother
% ---------- End .toc silence ----------
\hypersetup{
  colorlinks=true,
  linkcolor=accent,
  urlcolor=accent,
  citecolor=accent
}

% ---------- Minted (syntax highlighting) ----------
% NOTE: Compile with: -shell-escape
% ---------- Minted compatibility layer (CI-safe) ----------
% This document can render code with minted when -shell-escape is enabled, and
% falls back to listings when it is not (so CI builds do not hard-fail).
\newif\ifUseMinted
\UseMintedfalse
\begingroup
\ifdefined\pdfshellescape
  \ifnum\pdfshellescape=1\relax
    \global\UseMintedtrue
  \fi
\fi
\endgroup

\ifUseMinted
\else
\usepackage{fancyvrb} % provides \VerbatimEnvironment used by some wrappers

  % Provide a "listing" float compatible with minted's newfloat option
  \makeatletter
  \@ifundefined{c@listing}{%
    }{}
  \makeatother

  % Minimal language definitions for common "minted" lexers / labels used in docs.
  % These are intentionally lightweight; they exist primarily to avoid build failures.
  

